%%======================================================================
%%  Annexes
%%======================================================================
\chapter*{Annexes}
\addcontentsline{toc}{chapter}{Annexes}

%%----------------------------------------------------------------------
\section*{Annexe A : Liste complète des entités JPA}
\addcontentsline{toc}{section}{Annexe A : Liste complète des entités JPA}

\begin{longtable}{clll}
\caption{Liste des 25 entités du backend} \\
\toprule
\textbf{\#} & \textbf{Entité} & \textbf{Table} & \textbf{Type} \\
\midrule
\endfirsthead
\toprule
\textbf{\#} & \textbf{Entité} & \textbf{Table} & \textbf{Type} \\
\midrule
\endhead
\bottomrule
\endfoot
1  & User           & users             & Classe abstraite (base) \\
2  & Admin          & users             & Hérite de User \\
3  & Enseignant     & users             & Hérite de User \\
4  & Etudiant       & users             & Hérite de User \\
5  & Departement    & departements      & Entité indépendante \\
6  & Groupe         & groupes           & Entité indépendante \\
7  & Classe         & classes           & Entité indépendante \\
8  & Cours          & cours             & Entité avec relations \\
9  & Matiere        & matieres          & Entité avec relations \\
10 & Note           & notes             & Entité avec relations \\
11 & Devoir         & devoirs           & Entité avec relations \\
12 & Soumission     & soumissions       & Entité avec relations \\
13 & Evaluation     & evaluations       & Entité avec relations \\
14 & Absence        & absences          & Entité avec relations \\
15 & Justification  & justifications    & Entité avec relations \\
16 & EmploiDuTemps  & emploi\_du\_temps & Entité avec relations \\
17 & Inscription    & inscriptions      & Entité avec relations \\
18 & Projet         & projets           & Entité avec relations \\
19 & Demande        & demandes          & Entité avec relations \\
20 & Demande\_Stage & demandes\_stage   & Entité avec relations \\
21 & Reclamation    & reclamations      & Entité avec relations \\
22 & ReclamationNote& reclamation\_notes& Entité avec workflow \\
23 & Notification   & notifications     & Entité avec workflow \\
24 & Commentaire    & commentaires      & Entité (soft delete) \\
25 & Role           & —                 & Enum (ADMIN, ENSEIGNANT, ETUDIANT) \\
\end{longtable}

%%----------------------------------------------------------------------
\section*{Annexe B : Liste complète des contrôleurs REST}
\addcontentsline{toc}{section}{Annexe B : Liste complète des contrôleurs REST}

\begin{longtable}{cllc}
\caption{Les 28 contrôleurs REST du backend} \\
\toprule
\textbf{\#} & \textbf{Contrôleur} & \textbf{URL de base} & \textbf{Endpoints} \\
\midrule
\endfirsthead
\toprule
\textbf{\#} & \textbf{Contrôleur} & \textbf{URL de base} & \textbf{Endpoints} \\
\midrule
\endhead
\bottomrule
\endfoot
1  & HelloController           & /api            & 2 \\
2  & UserController            & /api/users      & 9 \\
3  & AdminController           & /api/admins     & 6 \\
4  & EnseignantController      & /api/enseignants& 7 \\
5  & EtudiantController        & /api/etudiants  & 7 \\
6  & DepartementController     & /api/departements& 6 \\
7  & GroupeController          & /api/groupes    & 6 \\
8  & ClasseController          & /api/classes    & 7 \\
9  & CoursController           & /api/cours      & 9 \\
10 & MatiereController         & /api/matieres   & 5 \\
11 & NoteController            & /api/notes      & 7 \\
12 & DevoirController          & /api/devoirs    & 7 \\
13 & SoumissionController      & /api/soumissions& 9 \\
14 & AbsenceController         & /api/absences   & 7 \\
15 & JustificationController   & /api/justifications & 10 \\
16 & EmploiDuTempsController   & /api/emploi-du-temps& 6 \\
17 & InscriptionController     & /api/inscriptions& 8 \\
18 & ProjetController          & /api/projets    & 6 \\
19 & EvaluationController      & /api/evaluations& — \\
20 & DemandeController         & /api/demandes   & 6 \\
21 & Demande\_StageController  & /api/demandes-stage& 6 \\
22 & ReclamationController     & /api/reclamations& 6 \\
23 & ReclamationNoteController & /api/reclamation-notes& 8 \\
24 & NotificationController    & /api/notifications& 14 \\
25 & DashboardController       & /api/dashboard  & 3 \\
26 & ActionsController         & /api/actions    & 4 \\
27 & CommentaireController     & /api/commentaires& 5 \\
28 & FileDownloadController    & /api/files      & 2 \\
\end{longtable}

%%----------------------------------------------------------------------
\section*{Annexe C : Schéma de la base de données}
\addcontentsline{toc}{section}{Annexe C : Schéma de la base de données}

\begin{table}[H]
\centering
\caption{Tables de la base de données}
\begin{tabularx}{\textwidth}{lXc}
\toprule
\textbf{Table} & \textbf{Description} & \textbf{Colonnes} \\
\midrule
users & Utilisateurs (héritage SINGLE\_TABLE) & 14 \\
departements & Départements & 5 \\
groupes & Groupes d'étudiants & 6 \\
classes & Classes d'enseignement & 6 \\
cours & Cours/Modules & 9 \\
cours\_etudiants & Table de jointure Cours ↔ Étudiant & 2 \\
matieres & Matières & 7 \\
notes & Notes des étudiants & 8 \\
devoirs & Devoirs/Travaux assignés & 8 \\
soumissions & Soumissions de devoirs & 11 \\
evaluations & Évaluations planifiées & 9 \\
absences & Absences enregistrées & 8 \\
justifications & Justifications d'absences & 8 \\
emploi\_du\_temps & Emploi du temps & 7 \\
inscriptions & Inscriptions étudiant-classe & 8 \\
projets & Projets collaboratifs & 7 \\
projet\_etudiants & Table de jointure Projet ↔ Étudiant & 2 \\
demandes & Demandes administratives & 9 \\
demandes\_stage & Demandes de stage & 8 \\
reclamations & Réclamations & 8 \\
reclamation\_notes & Réclamations de notes & 12 \\
notifications & Notifications & 13 \\
commentaires & Commentaires & 7 \\
\bottomrule
\end{tabularx}
\end{table}

%%----------------------------------------------------------------------
\section*{Annexe D : Collection Postman}
\addcontentsline{toc}{section}{Annexe D : Collection Postman}

Une collection Postman complète est fournie avec le projet dans le fichier :

\texttt{docs/Education\_Platform\_API.postman\_collection.json}

\vspace{0.3cm}

Cette collection contient l'ensemble des requêtes HTTP préconfigurées pour tester les 90+ endpoints de l'API, organisées par ressource. Elle peut être importée directement dans Postman.

\vspace{0.3cm}

\textbf{Pour l'importer :}
\begin{enumerate}[noitemsep]
    \item Ouvrir Postman ;
    \item Cliquer sur « Import » ;
    \item Sélectionner le fichier JSON ;
    \item Toutes les requêtes sont alors disponibles avec les URL et les corps pré-remplis.
\end{enumerate}

\newpage
